\documentclass[a4paper]{article}
% Add "final" option in the article to remove todo
\usepackage[common_math_textnormal,todonotes,math_base_note,math_simple]{gatmeo}
\usepackage{tikz-cd}
\newcommand{\mynewpage}{\newpage}
\fancyhead[R]{\thepage}


\title{\bf{
                Kazhdan-Lusztig-Stanley polynomials and Chow polynomials for Bia{\l}ynicki-Birula stratified spaces
}}

\newcommand\Gm{{\mathbb{G}_m}}
\newcommand\Fq{{\mathbb{F}_q}}
\newcommand\Ql{{\mathbb{Q}_l}}
\newcommand\LL{{\mathbb{L}}}
\newcommand\II{{\mathcal{I}}}
\newcommand\Qlb{{\overline{\mathbb{Q}}_l}}
\newcommand\Fqs{{\mathbb{F}_{q^s}}}
\newcommand\Fqb{{\overline{\mathbb{F}}_q}}
\newcommandx\IP{\mathcal{I}(P)}
\newcommandx*\CHM[1][1=M]{\text{CH}(#1)}
\newcommandx*\XE[1][1=E]{X_{#1}}
\newcommandx*\aXE[1][1=E]{\underline{X}_{#1}}
\newcommandx*\XV[1][1=V]{X_{#1}}
\newcommandx*\aXV[1][1=V]{\underline{X}_{#1}}
\newcommandx*\aCHM[1][1=M]{\underline{\text{CH}}(#1)}
\newcommandx\Bl[2]{\text{Bl}_{#1}(#2)}

\begin{document}
\maketitle
\vspace{-3.5em}
%\input{sections/abstract}

\thispagestyle{empty}
\tableofcontents
\listoftodos

\section{Introduction}

Kazhdan-Lusztig-Stanley (KLS) polynomials and Chow polynomials are elements in the incidence algebra of a (locally finite) weakly ranked poset that satisfy some recurrence relations encoded by the poset. 
We are interested in interpreting the polynomials geometrically.

Let $X$ be a variety over $\mathbb{F}_q$ with a smooth stratification 
\[
        X=\bigsqcup _{x\in P}V_x,
\]
i.e. each stratum $V_x$ is a smooth subvariety of $X$ and the closure of each stratum is a union of strata. We define a partial order on $P$ by putting $x\leq y$ if and only if $V_x\subseteq \overline{V_y}$, and a weak rank function by $r_{xy}:=\dim V_y - \dim V_x$.
Proudfoot~\cite{Proudfoot2018} introduces a geometric framework in which point counts define a $P$-kernel and intersection cohomology realizes the associated KLS functions. He shows that all Schubert varieties and toric varieties give rise to KLS-polynomials.

In the case of $X$ being a toric variety, Micha{\l}ek, Monin, and Wang~\cite{MichalekMoninWang2026} consider a different (possibly singular) stratification, the Bia{\l}ynicki-Birula (BB) decomposition. The same framework for constructing an element in the incidence algebra extends. They show that if the Chow quotient of $X$ is rationally smooth, then the entries of the constructed element are KLS-polynomials.
They show that the constructed Chow polynomials are the Poincar\'e polynomials of the Chow quotient.

In this paper, we unify the two theories. In the case of $X$ being a $\Gm$-variety, we give sufficient conditions on $X$ such that we can construct KLS-polynomials. We also give sufficient conditions on the Chow quotient of $X$ such that we can construct Chow polynomials.

\section{Combinatorics}

A \vocab{poset} is a set $P$ with a reflexive, antisymmetric, and transitive relation $\leq$. We say that $P$ is locally finite if, for all $x\leq y\in P$,  the set
\[
        [x,y]:=\{z\in P\mid x\leq z\leq y\}
\]
is finite. 
A \vocab{weak rank function} assigns an integer $r_{xy}$ to each interval, with $r_{xx}=0$, $r_{xy}>0$ for $x<y$, and
\[
        r_{xz}=r_{xy}+r_{yz}\qquad(x\leq y\leq z).
\]
We call $P$ together with such a function a \vocab{weakly ranked poset}.

For a commutative ring $R$ with identity, define its \vocab{incidence algebra} by
\[
        \II(P,R):=\{F:\{(x,y)\in P^2\mid x\leq y\}\longrightarrow R\}.
\]
If $P$ is locally finite, $\II(P,R)$ admits a ring structure with pointwise addition and product given by convolution:
\[
        (FG)_{xz}=\sum_{x\leq y\leq z}F_{xy}G_{yz}.
\]
The identity is the element $\delta$ with $\delta_{xx}=1_R$ and $\delta_{xy}=0$ for $x<y$. When the coefficient ring is omitted, we use the convention
\[
        \II(P):=\II(P,\mathbb{Z}[t]).
\]

\begin{lemma}[lem:incidence_inverse]{}
        An element $F\in\II(P,R)$ has a two-sided inverse if and only if $F_{xx}\in R^\times$ for every $x\in P$. 
        \begin{proof}
                The proof of \cite[Lemma~2.1]{Proudfoot2018} extends for any commutative ring $R$.
        \end{proof}
\end{lemma}

For a weakly ranked poset $P$, let
\[
        \II_r(P) :=
        \left\{ F\in\II(P) \,\middle|\, \deg F_{xy}\leq r_{xy} \text{ for every }x\leq y \right\}.
\]
The additivity of the weak rank function implies that $\II_r(P)$ is a subalgebra of $\II(P)$. 
Define the \vocab{bar involution} on $\II_r(P)$ by
\[
        \overline{F}_{xy}(t)
        :=
        t^{r_{xy}}F_{xy}(t^{-1}).
\]
The weak rank identity $r_{xz}=r_{xy}+r_{yz}$ implies $\overline{FG}=\overline F\,\overline G$ and $\overline{\overline F}=F$.

\begin{definition}{}
        An element $\kappa\in\II_r(P)$ is called a \vocab{$P$-kernel} if $\kappa_{xx}=1$ for every $x\in P$ and 
        \[
                \overline{\kappa}\,\kappa=\delta.
        \]
\end{definition}

\section{Geometric Setup}

We follow the notation of Proudfoot~\cite[Section~3]{Proudfoot2018}. In our case, some of his assumptions are automatic.

Fix a finite field $\mathbb{F}_q$ with algebraic closure $\overline{\mathbb{F}_q}$. Let $l$ be a prime not dividing $q$. For any variety $Z$ over $\mathbb{F}_q$, denote by $IC_Z$ the $l$-adic intersection cohomology sheaf on the variety $Z(\overline{\mathbb{F}}_q)$ with the convention that we do not shift $IC_Z$ to make it perverse. In particular, if $Z$ is smooth, $IC_Z$ is isomorphic to the constant sheaf in degree zero.

Let $X$ be a normal projective $\Gm$-variety over $\mathbb{F}_q$ with finitely many isolated fixed points $P:=X^\Gm$. After replacing $\mathbb{F}_q$ by a finite extension if necessary, we assume that every fixed point is $\mathbb{F}_q$-rational.

The positive and negative Bia{\l}ynicki-Birula cells of $X$ corresponding to a fixed point $x\in X^\Gm$ are defined to be
\[
        X^+_x=\CB{p\in X\mid \lim_{t\to 0}t\cdot p=x}\textnormal{ and }
        X^-_x=\CB{p\in X\mid \lim_{t\to \infty }t\cdot p=x}
\]
respectively. $X^+_x$ and $X^-_x$ can be singular in general. 
For each $x,y\in P$, define
\[
        U_{xy}:=\left(X^-_x\times _X X^+_y\right)_{\textnormal{red}}\textnormal{ and }
        Z_{xy}:=\left(X^-_x\times _X \overline{X^+_y}\right)_{\textnormal{red}}.
\]

Assume further that the BB-decomposition 
\[
        X=\bigsqcup _{x\in P}X^+_x
\]
is a stratification, i.e. the closure of each stratum is a union of strata. We define a partial order on $P$ by putting $x\leq y$ if and only if $X^+_x\subseteq \overline{X^+_y}$, and a weak rank function by $r_{xy}:=\dim X^+_y - \dim X^+_x$.
Hence 
\[
        \overline{X^+_y}=\bigsqcup_{x\leq y}X_x^+.
\]
Geometrically, a point in $U_{xy}$ gives a $\Gm$-orbit whose $t\rightarrow 0$ endpoint is $y$ and whose $t\rightarrow \infty $ endpoint is $x$, i.e. $x\in \overline{X^+_y}$.
In particular, $U_{xy}\neq \emptyset$ implies $x\leq y$.
Hence
\begin{equation}
        Z_{xz}=\bigsqcup_{x\leq y\leq z}U_{xy}.
        \label{eq:Z_decompose_as_U}
\end{equation}

By Drinfeld's theory of attractors and repellers, the limit morphism
from the attractor to the fixed-point locus is affine
\cite[Theorem~1.4.2]{Drinfeld2013}. Applying this to the inverse
$\Gm$-action, and using that $x$ is an isolated fixed point, we see
that $X_x^-$ is affine. Moreover, the inverse $\Gm$-action on $X_x^-$
extends to an $\mathbb A^1$-action contracting $X_x^-$ to $x$.
Consequently,
\[
        \mathbb F_q[X_x^-]=\bigoplus_{i\geq 0}A_i,
        \qquad A_0=\mathbb F_q.
\]
Since $\overline{X_z^+}$ is closed and $\Gm$-invariant,
\[
        Z_{xz}=X_x^-\times_X\overline{X_z^+}
\]
is defined in $X_x^-$ by a homogeneous ideal. Hence $Z_{xz}$ is a
weighted affine cone with cone point $x$.

We write $IH_p^*(Z)$ to denote the cohomology of the stalk of $IC_Z$ at $p\in Z$. For $x,y\in P$, we denote $IH^*_{xy}:=IH^*_{x}(\overline{X^+_y})$. For a weighted affine cone, contraction to the cone point induces a Frobenius-equivariant isomorphism~\cite[Lemma~4.5(a)]{KazhdanLusztig1980} 

\begin{equation}
        IH_x(Z_{xz})=IH(Z_{xz}).
        \label{eq:IH_cone_pt}
\end{equation}

Let $W_{\mathbb F_q}\cong\mathbb Z$ be the Weil group of $\mathbb F_q$, generated by geometric Frobenius $Fr_q$, and put
\[
        R_q:=K_0\!\left(\operatorname{Rep}_{\Qlb}(W_{\mathbb F_q})\right).
\]
Here $\operatorname{Rep}_{\Qlb}(W_{\mathbb F_q})$ is the category of finite-dimensional $\Qlb$-representations of $W_{\mathbb F_q}$, equivalently finite-dimensional $\Qlb$-vector spaces with an invertible endomorphism. For such a representation $V$, write $[V]\in R_q$ for its class. Addition is induced by short exact sequences and multiplication by tensor product. Traces extend additively to virtual classes in $R_q$.

\begin{lemma}[lem:frobenius_traces]{}
        For $A,B\in R_q$, we have $A=B$ if and only if
        \[
                \operatorname{Tr}(Fr_q^s;A)=\operatorname{Tr}(Fr_q^s;B)
                \qquad\text{for every }s\geq 1.
        \]
        \begin{proof}
                Equality of classes implies equality of traces. Conversely, let $\alpha_1,\dots,\alpha_m$ be the distinct Frobenius eigenvalues in $A-B$, with virtual multiplicities $n_1,\dots,n_m$. Equality of traces gives
                \[
                        \sum_{i=1}^m n_i\alpha_i^s=0\qquad(s\geq 1).
                \]
                For $s=1,\dots,m$, the coefficient matrix is a Vandermonde matrix multiplied by the invertible diagonal matrix with entries $\alpha_1,\dots,\alpha_m$. Thus all $n_i$ vanish, $A=B$.
        \end{proof}
\end{lemma}
The ring $R_q$ remembers the eigenvalues with multiplicities but forgets Jordan blocks, just as traces of Frobenius powers do. In contrast, $K_0(\operatorname{Vect}_{\Qlb})\cong\mathbb Z$ remembers only dimension.

Duality defines an involution
\[
        [V]^\vee:=[V^\vee],
\]
where Frobenius acts on $V^\vee$ by the inverse transpose. For a bounded complex $C^\bullet$ with Frobenius action, put
\[
        [C^\bullet]:=\sum_i(-1)^i[H^i(C^\bullet)]\in R_q.
\]
Let
\[
        \LL:=[\Qlb(-1)]\in R^\times_q.
\]
Geometric Frobenius acts on this Tate twist by $q$, so
\[
        \operatorname{Tr}(Fr_q^s;\LL^j)=q^{js}
        \qquad(s\geq 1,\ j\in\mathbb Z).
\]

The following theorems play a central role in our proof:
\begin{theorem}[thm:Grothendieck_Lefschetz_formula]{Grothendieck--Lefschetz trace formula \cite[III.12.1(4)]{KiehlWeissauer2001}}
        Let $Z$ be separated and of finite type over $\Fq$, and let $\mathcal F$ be a constructible $l$-adic complex with Weil structure. For every $s\geq 1$,
        \[
                \operatorname{Tr}\!\left(Fr_q^s;[R\Gamma_c(Z_{\Fqb},\mathcal F)]\right)
                =\sum_{u\in Z(\Fqs)}
                \operatorname{Tr}\!\left(Fr_{q^s,u};[\mathcal F_{\bar u}]\right),
        \]
        where $\bar u$ is a geometric point above $u$ and $Fr_{q^s,u}$ denotes geometric Frobenius acting on the stalk at $\bar u$. In particular,
        \[
                \operatorname{Tr}\!\left(Fr_q^s;[R\Gamma_c(Z_{\Fqb},\Qlb)]\right)
                =|Z(\Fqs)|.
        \]
\end{theorem}

\begin{theorem}[thm:Verdier_duality]{Verdier duality \cite[II.7.3]{KiehlWeissauer2001}}
        Let $a:Z\to\operatorname{Spec}\Fq$ be separated and of finite type. For a constructible $l$-adic complex $\mathcal F$, there is a canonical isomorphism
        \[
                R\Gamma_c(Z_{\Fqb},\mathcal F)^\vee
                \simeq R\Gamma(Z_{\Fqb},\mathbb D_Z\mathcal F),
        \]
        where $\mathbb D_Z\mathcal F:=R\mathcal{H}om(\mathcal F,a^!\Qlb)$ is the Verdier dual. If $\mathcal F$ carries a Weil structure, this isomorphism is Frobenius-equivariant.

        If $Z$ is pure of dimension $d$, then, with our unshifted normalization,
        \[
                \mathbb D_Z(IC_Z)\simeq IC_Z(d)[2d].
        \]
        Consequently,
        \[
                [R\Gamma_c(Z_{\Fqb},IC_Z)]
                =\LL^d[R\Gamma(Z_{\Fqb},IC_Z)]^\vee
                \qquad\text{in }R_q.
        \]
\end{theorem}

\begin{definition}[def:geometric_hypotheses]{}
        We say that $X$ satisfies \vocab{Frobenius-compatible IC-locality} if, for every $x\leq y\leq z$, every $s\geq 1$, and every $u\in U_{xy}(\Fqs)$, we have an isomorphism of graded Frobenius modules
        \[
                \left(IH_u^\bullet(Z_{xz}),Fr_{q^s,u}\right)
                \simeq
                \left(IH_{yz}^\bullet,Fr_q^s\right).
        \]
        We say that $X$ satisfies \vocab{polynomial point-count} if there is $\kappa\in\II_r(P)$ such that, for every $x\leq y$ and every $s\geq 1$,
        \[
                |U_{xy}(\Fqs)|=\kappa_{xy}(q^s).
        \]
        We say that $X$ is \vocab{dimensionally transverse} if for every $x\leq z$,
        \[
                Z_{xz}\text{ is pure of dimension }r_{xz}.
        \]
\end{definition}

For $x\leq z$, put
\[
        F_{xz}:=[IH_{xz}^\bullet]
        =\sum_i(-1)^i[IH_{xz}^i]\in R_q.
\]
Then $F\in\II(P,R_q)$. Define the bar operation on $\II(P,R_q)$ by
\[
        \overline{F}_{xz}:=\LL^{r_{xz}}F_{xz}^\vee.
\]
Since $r_{xz}=r_{xy}+r_{yz}$ and $(\LL^aV)^\vee=\LL^{-a}V^\vee$, this is an algebra involution.
Define evaluation at $\LL$ entrywise by
\[
        \kappa(\LL)_{xz}:=\kappa_{xz}(\LL).
\]
Since $\overline{\kappa(\LL)}_{xz}=\LL^{r_{xz}}\kappa_{xz}(\LL^{-1})$, 
evaluation is compatible with the bar operation, i.e. $\overline{\kappa(\LL)}=\overline{\kappa}(\LL)$.

\begin{theorem}[thm:]{}
        Let $X$ be as above, satisfying Frobenius-compatible IC-locality and polynomial point-count, and assume that $X$ is dimensionally transverse. Then $\kappa$ is a $P$-kernel.
        \begin{proof}
                Take $x\leq z$. By Frobenius-compatible IC-locality with $y=x,u=x$, and $s=1$, and \Cref{eq:IH_cone_pt},
                \[
                        F_{xz}=[IH_{xz}^\bullet]=[IH_x^\bullet(Z_{xz})]=[IH^\bullet(Z_{xz})].
                \]
                Since $Z_{xz}$ is pure of dimension $r_{xz}$, \Cref{thm:Verdier_duality} gives
                \[
                        [IH_c^\bullet(Z_{xz})]=\LL^{r_{xz}}F_{xz}^\vee=\overline{F}_{xz}.
                \]
                Now fix $s\geq 1$. By \Cref{thm:Grothendieck_Lefschetz_formula}, \Cref{eq:Z_decompose_as_U}, Frobenius-compatible IC-locality, and polynomial point-count,
                \begin{align*}
                        \operatorname{Tr}(Fr_q^s;\overline{F}_{xz})
                        &=\operatorname{Tr}(Fr_q^s;[IH_c^\bullet(Z_{xz})])\\
                        &=\sum_{x\leq y\leq z}\sum_{u\in U_{xy}(\Fqs)}
                        \operatorname{Tr}\!\left(Fr_{q^s,u};[IH_u^\bullet(Z_{xz})]\right)\\
                        &=\sum_{x\leq y\leq z}|U_{xy}(\Fqs)|\operatorname{Tr}(Fr_q^s;F_{yz})\\
                        &=\sum_{x\leq y\leq z}\kappa_{xy}(q^s)\operatorname{Tr}(Fr_q^s;F_{yz})\\
                        &=\operatorname{Tr}\!\left(Fr_q^s;(\kappa(\LL)F)_{xz}\right).
                \end{align*}
                Since this holds for every $s\geq 1$, \Cref{lem:frobenius_traces} implies
                \[
                        \overline{F}=\kappa(\LL)F.
                \]
                Applying the bar involution gives
                \[
                        F=\overline{\kappa}(\LL)\overline{F}
                        =\overline{\kappa}(\LL)\kappa(\LL)F.
                \]
                Since $Z_{xx}$ is a zero-dimensional weighted affine cone with cone point $x$, $Z_{xx}=\{x\}$. Hence $F_{xx}=1_{R_q}$, $F$ is invertible by \Cref{lem:incidence_inverse}. Cancelling $F$ gives
                \[
                        \overline{\kappa}(\LL)\kappa(\LL)=\delta.
                \]
                The proof follows from the homomorphism
                \[
                        \mathbb Z[t,t^{-1}]\longrightarrow R_q,\qquad t\longmapsto\LL,
                \]
                is injective, and $U_{xx}=Z_{xx}$ has one rational point over every finite extension.
        \end{proof}
\end{theorem}

\begin{remark}
        The theorem does not require chastity.
        If, in addition, every $IH_{pr}^\bullet$ is chaste, i.e. $IH_{pr}^{2i+1}=0$ and geometric Frobenius acts on $IH_{pr}^{2i}$ by multiplication by $q^i$, then
        \[
                F_{xz}=f_{xz}(\LL),
                \qquad
                f_{xz}(t):=
                \sum_{i\geq0}\dim IH_{xz}^{2i}\,t^i.
        \]
        Indeed, $Tr(Fr_q^s;F_{xz}) = f_{xz}(q^s)$.  Under this identification, the involution
        \[
                F_{xz}\longmapsto
                \LL^{r_{xz}}F_{xz}^\vee
        \]
        becomes the usual bar involution
        \[
                f_{xz}(t)\longmapsto
                t^{r_{xz}}f_{xz}(t^{-1}).
        \]
        Thus, under chastity, the argument above reduces to Proudfoot's polynomial-valued proof \cite[Section~3.2]{Proudfoot2018}.
\end{remark}

%Bialynichi-Birula original work is on smooth complete $\Gm$-variety. He shows that the cells are locally closed subvariety isomorphic to some affine spaces $\mathbb{A}^n$. 

\section{Peterson Variety}

Let $Y\in \mathfrak{g}$, $H\subseteq \mathfrak{g}$ s.t. $[H,\mathfrak{b}]\subseteq H$. Define
\[
  Hess(Y,H):=\CB{gB:Ad(g^{-1})Y\in H}\subseteq G/B.
\]
In type A, $H$ can be described using a non-decreasing function $h:[n]\rightarrow [n]$ such that $h(i)\geq i$. Pick a standard flag $(E_1,E_2,\dots ,E_n)\in Fl(\mathbb{C}^n)$. Define
\[
  H_h:=\CB{A\in \mathfrak{gl}_n\mid AE_i\subseteq E_{h_i}}.
\]
An alternative describsion of Hessenberg variety is
\[
  Hess(X,H_h) =\CB{V_\bullet\in Fl(\mathbb{C}^n)\mid XV_j\subseteq V_{h_j}\ \forall\ 1\leq j\leq n}.
\]

We focuse on regular nilpotent Hessenberg variety, i.e. when $Y=N:=E_{12}+E_{23}+\cdots +E_{n-1n}$. In this case, we have a $\Gm$ action given by
\[
  \lambda(t)=diag(t,t^2,\dots ,t^n).
\]
This can be seen from the fact that $\lambda(t)^{-1}N\lambda(t)=tN$.
\begin{definition}[def:]{}
  Peterson variety is defined as $Hess(N,H_{Pet})$ with $Pet(i)=\min(i+1,n)$.
\end{definition}

Since $\lambda$ is regular, the fixed points $Hess(N,H_h)^{\lambda(\mathbb{C}^\times )}$ are the permutation flags $wB$.
One way to study Peterson variety is to look at its intersection with the bruhat cells $BwB/B$. Since $w^{-1}Nw=\sum_{i=1}^{n-1}E_{w^{-1}(i), w^{-1}(i+1)}$, it lies in $H_{Pet}$ if and only if either $w^{-1}(i)>w^{-1}(i+1)$, or $w^{-1}(i)=w(i+1)+1$, or equivalently $w=w_J$ be the longest element of $W_J=\AB{s_i\mid i\in J }$ for some $J\subseteq [n-1]$.
Indeed, we have the following:

\begin{lemma}[lem:]{}
  The $\mathbb{C}^\times $ fixed points are of Peterson variety are $\CB{w_J\mid J\subseteq [n-1]}$. Furthermore, the BB-cells are precisely the intersection with the Bruhat cells, i.e.
  $X_J^+=Pet_n\cap BwJB/B$.
  \begin{proof}
    For $x\in BwB/B$, $x=bwB$. Write $b=ut$, we can absorb $t$ on the right as $w^{-1}tw\in T\subseteq B$, $twB=w(w^{-1}tw)B=wB$. Hence we can write $x=uwB$. Since $\lim_{t\rightarrow 0}\lambda(t)^{-1}u\lambda(t)\rightarrow 1$, we have
    \[
      \lim_{t\rightarrow 0}\lambda(t)\cdot uwB=\lim_{t\rightarrow 0}(\lambda(t)^{-1}u\lambda(t))wB=wB.
    \]
  \end{proof}
\end{lemma}

We would like to give a description of the positive BB-cells:
\begin{lemma}[lem:coordinate_XJ]{}
  Write $J\subseteq [n-1]$ as disjoint union of connected commponents $J=C_1\sqcup\cdots \sqcup C_s$, where $C_i=[p_i,q_i]$ and $p_{i+1}-q_i>1$. Define
  \[
    N_i=E_{p_i,p_i+1}+\cdots +E_{q_i,q_i+1}.
  \]
  Then
  \[
    X_J^+=\CB{u_J(a)w_JB; a_{i,j}\in \mathbb{C}},\quad u_J(a):=\prod_{i}^{}\left(1+a_{i,1}N_{i}+\cdots +a_{i,|C_i|}N_{i}^{|C_i|}\right).
  \]
  In particular, $X^+_J\simeq \mathbb{A}^{|J|}$. For $J=[n-1]$, we have $X_{[n-1]}^+=Z_U(N)$.
  \begin{proof}
    We prove the case of $J=[n-1]$. The general case follow from breaking down the calculation block by block labelled by $C_i$. Write $x\in Bw_0B/B$ as $uw_0B$ as above. Peterson condition says $w_0^{-1}u^{-1}Nuw_0\in H_{Pet}$. Note that $w_0H_{Pet}w_0^{-1}=\mathfrak{b}^-+\CB{\textnormal{positive simple root spaces}}$. But $u^{-1}Nu\in \mathfrak{n}^+=Lie(U)$. This forces $u^{-1}Nu=N$. The proof concludes by noticing $Z_U(N)=\CB{1+a_1N+a_2N;=2+\cdots +a_{n-1}N^{n-1}}$.
  \end{proof}
\end{lemma}
\begin{remark}
  Since $\lambda(t)^{-1}N\lambda(t)=tN$, $\lambda^{-1}N^s\lambda(t)=t^sN$. $\lambda$ acting on $u_J(a)$ with weights $t^s$ on coordinate $a_{i,s}$.
\end{remark}

Next we are interested in studying $U_{I,J}=X_I^{-}\cap X_J^+$. In particular, we are asking for $x=MB\in X_J^+$, when is $x\in B^{-1}w_IB/B$. Let $V_k=\Span\CB{\textnormal{first }k\textnormal{ columns of }M}$, and let the opposite coordinate flag be $E_r^-=\Span(e_{n-r+1,\dots ,e_n})$. Then $x\in B^-w_IB/B$ iff
\[
  \dim(V_k\cap E_r^-)=\#\CB{j\leq k\mid w_I(j)\geq n-r+1}.
\]

\begin{lemma}[lem:]{}
  Let $x=u_J(a)w_JB\in B_+w_JB/B$. Denote by $\Delta_i$ the $i$-th principal minor of $u_J(a)w_J$. Then $x\in B^-w_IB/B$ iff $\Delta_i=0$ for $i\in I$ and $\Delta_i\neq 0$ if $i\notin I$.
  \begin{proof}
    Observe that $\Delta_k(lwb)=0$ iff $\Delta_k(w)=0$.
    The proof conlcudes by $\Delta_k(w)\neq 0$ iff $w([k])=[k]$. More precisely, by construction, for $k\neq I$, we have $w_I([k])=[k]$, and for $k\in I$, we have $w_I([k])\neq [k]$.
  \end{proof}
\end{lemma}

\begin{example}[exp:]{}
  For $Pet_3$, $M=\pmx{b&a&1\\a&1&0\\1&0&0}$. The two principal minors are $\Delta_1=b,\Delta_2=b-a^2$. $U_{\emptyset ,12}=\CB{a,b\mid b\neq 0,b-a^2\neq 0}\simeq \Gm^2$. 

  $U_{I,J}$ can be irreducible in general.
  For $Pet_4$, the 3 principal minors are $\Delta_1=c$, $\Delta_2=ac-b^2$, $\Delta_3=-a^3+2ab-c$. For $I=\CB{1,3}$, $J=\CB{1,2,3}$, we have $c=0, -b^2\neq 0$, and $a(2b-a^2)=0$. We see that $U_{13,123}\simeq \Gm\sqcup \Gm$.

  For $I=\CB{1}$, $J=\CB{1,2,3}$, we instead of $c=0,b\neq 0,a\neq 0,2b-a^2\neq 0$. Note that $b/a^2$ is $\Gm$ invariant by the remark preceeding \Cref{lem:coordinate_XJ}. Hence 
  \[
    Q_{1,123}:=U_{1,123}/\Gm=\mathbb{A}^1\setminus \CB{0,1/2}\simeq \mathbb{P}^1\setminus \CB{0,1/2,\infty }.
  \]
  We expect the boundary of the compactification $Q_{1,123}$ correspond to the broken trajectories from fixed point $1$ to $123$.
  We can see that $0$ correspond to $b=0$, hence $\Delta_2=0$ as $c=0$. This coorespond to $U_{12,123}$.
  On the other hand, $\frac{1}{2}$ and $\infty $ corresponds to the two $\Gm$ compnents in $U_{13,123}$ respectively.
\end{example}

\begin{example}[exp:]{}
  Point counting for $Pet_6$ fails to give P-kernel. Consider $I=\CB{1,3,5}$, $J=\CB{1,2,3,4,5}$. $U_{I,J}$ is given by $\Delta_1=\Delta_3=\Delta_5=0$ and $\Delta_2\Delta_4\neq 0$. Using coordinates $(a,b,c,d,e)$ for $X_J^+$, 
\end{example}

<++>

\begin{remark}
  In the case of full flag variety, there is a automorphism exchanging $X^+$ and $X^-$ given by
  \[
    L_{w_0}:G/B\rightarrow G/B, gB\mapsto w_0gB.
  \]
  Since $w_0Bw_0^{-1}=B^-$, we get
  \[
    L_{w_0}(BwB/B)=B^-(w_0w)B/B.
  \]
  Notice that $w_0\lambda(t)w_0^{-1}=t^{n+1}\lambda(t)^{-1}$. Since $t^{n+1}I$ acts trivially on $G/B$, $L_{w_0}$ reverse the flow of the BB cells.
  But this automorphism does not preserve Peterson variety, it sends Peterson to a different Hessenberg variety:
  \[
    L_{w_0}(Hess(N,H_{Pet}))=Hess(w_0Nw_0^{-1},H_{Pet}).
  \]
  The difficulty lies in finding the intersections of the opposite cells with Peterson variety.
\end{remark}


\bibliographystyle{alpha}
\bibliography{main_kls}

\end{document}
